% Integration-ready contribution. Requires amsmath, amssymb, amsthm.
\providecommand{\shuffle}{\mathbin{\amalg}}
\section{The complete Cayley ideal does not close the bounded weight-seven search}
\label{gf:sec}

The pinned manuscript leaves the mixed Gaussian identities for $S_6$ and
$S_8$ conjectural. Two recent source contributions sharpened the algebraic
problem in complementary ways. One supplied an exact integer functional
separating the $S_6$ target from a specified collection of $5131$ relations
of depth at most four \cite{gf:Bounded}. The other constructed a
depth-preserving retraction for the \emph{full} shuffle ideal generated by
convergent Cayley identities \cite{gf:Depth}. It explicitly proposed
projecting the bounded rows by that retraction as a next proof attempt.

We complete that proposed calculation. The inherited integer functional
also annihilates the intersection of the \emph{entire} Cayley shuffle ideal
with the weight-seven space of depth at most four. This includes all
Cayley rows and all their shuffle multiples, even when their construction
uses higher-depth intermediate words. Thus completing the Cayley side
alone cannot prove $S_6$ from the stated bounded additional relations.
The numerical $S_6$ identity remains open. This is a result about an
explicit formal presentation; it is not nonvanishing of its evaluated
period residual, and it gives no period-independence assertion.

\subsection{The target and the exact formal presentation}

Use the convention
\[
 \operatorname{Li}_{r,s}(u,v)=\sum_{n>m\geq1}\frac{u^n v^m}{n^r m^s},
 \qquad g_{r,s}=\operatorname{Im}\operatorname{Li}_{r,s}(i,1),
 \qquad S_p=\sum_{n\geq0}\frac{(-1)^nH_n}{(2n+1)^p}.
\]
Let $L=\log2$, and write the conjectural residual in its already reconciled
short coordinate form:
\begin{align}
 R_6={}&S_6+\frac{12334}{527}g_{6,1}+\frac{384}{31}g_{5,2}
       +\frac{2136}{527}g_{4,3}-\frac{3179\pi^7}{5713920}\notag\\
      &+\frac{801}{2108}\beta(4)\zeta(3)+2\beta(6)L.
 \label{gf:eq:residual}
\end{align}
The conjecture is $R_6=0$.

For formal words take the positive forms
\[
 z=\frac{dt}{t},\quad c=\frac{dt}{1-t},\quad
 a=\frac{i\,dt}{1-it},\quad b=-\frac{dt}{1+t},\quad
 \bar a=-\frac{i\,dt}{1+it}.
\]
Their iterated integrals $H(w)$ are written outer letter first on $[0,1]$.
An admissible nonempty word begins with a letter other than $c$ and ends
with a letter other than $z$. Let $\mathcal A$ be the graded rational
shuffle algebra on the admissible words and the empty word. Conjugation
$K$ interchanges $a,\bar a$ and fixes $z,c,b$. The superscript $-$ denotes
its negative eigenspace, and $F_D\mathcal A$ is the span of words with
at most $D$ nonzero letters. This is the usual series-depth filtration.

The map induced by $t\mapsto(1-t)/(1+t)$ is
\[
 C(v_1\cdots v_n)=\tau(v_n)\cdots\tau(v_1),\qquad
 \begin{array}{c|ccccc}
 v&z&c&a&b&\bar a\\ \hline
 \tau v&c-b&z+b&b-\bar a&b&b-a.
 \end{array}
\]
Here $\tau=-\phi^*$ incorporates the reversal of orientation. Every
transformed admissible word is admissible, and change of variables gives
$H(Cw)=H(w)$. Consequently the period map annihilates the graded,
$K$-stable shuffle ideal
\begin{equation}
 \mathcal J=\langle f-Cf:f\in\mathcal A\rangle_{\shuffle}.
 \label{gf:eq:ideal}
\end{equation}
No restriction on weight, depth, or multiplier is made in this definition.

In the following polynomial, juxtaposition is concatenation and $\star$
is the usual series stuffle product, transferred to words by cumulative
colors. Define
\begin{align}
 P_6={}&128588z^5a\bar a+3138084z^5aa
          +1592832z^4aza+521184z^3az^2a-216172z^6a\notag\\
       &+48861\bigl((z^3a)\star(z^2c)\bigr)
          -257176\bigl((z^5a)\star b\bigr),\qquad
 T=\frac{P_6-KP_6}{2}.
 \label{gf:eq:target}
\end{align}
This is the exact target used by the earlier bounded certificate.
It belongs to $F_2\mathcal A_7^-$. The elementary harmonic split gives
$S_6=g_{6,1}+\operatorname{Im}\operatorname{Li}_{6,1}(i,-1)$, while
$\beta(7)=61\pi^7/184320$. Using also $H(b)=-L$ gives
\begin{equation}
 H(T)=128588\,iR_6.
 \label{gf:eq:target-value}
\end{equation}
Both products in \eqref{gf:eq:target} use convergent inputs, so their
stuffle expansion requires no regularization.

For precision, $\mathscr R_4\subset F_4\mathcal A_7^-$ denotes the
\emph{same} inherited finite row family as in \cite{gf:Bounded}. Its
construction is retained in the delivered independent enumerator:
\begin{enumerate}
\item All weight-seven shuffle-minus-stuffle products with sum of input
      depths at most four. Inputs are admissible except for the permitted
      single factor $\operatorname{Li}_1(1)$; the standard one-divergence
      comparison cancels its divergent term.
\item All convergent level-two-to-level-four distribution rows and their
      products with admissible complementary-weight words, provided the
      sum of depths is at most four. Products are expanded by stuffle.
\item The already proved imaginary $S_2$ and $S_4$ rows multiplied by real
      parts of all admissible complementary-weight words of depth at most
      two, again expanded by stuffle.
\item The inherited balanced Cayley rows and their specified products.
      In the adapted alphabet below, the seed $u$ begins with a letter
      other than $x$, ends with a letter other than $z$, and satisfies
      $\min(n_z(u),n_x(u))\geq1$. A multiplier $v$ is allowed when
      $|u|-\min(n_z(u),n_x(u))+\operatorname{dep}(v)\leq4$ and
      $|u|+|v|=7$. Its product is expanded by stuffle.
\end{enumerate}
Imaginary projection, rational normalization, and duplicate removal give,
in this order, $3460$, $1394$, $40$, and $237$ new nonzero rows, respectively:
$5131$ rows in total. These counts specify a formal enumeration, not
dimensions of periods. The fourth family may be retained even though the
new assertion explicitly imposes the complete ideal \eqref{gf:eq:ideal}.

\begin{theorem}[Full-Cayley extension of the weight-seven obstruction]
\label{gf:thm:main}
With the formal target and row family just specified,
\begin{equation}
 \boxed{\qquad
 T\notin\mathcal J_7^-+\operatorname{span}_{\mathbb Q}\mathscr R_4.
 \qquad}
 \label{gf:eq:nonmembership}
\end{equation}
In particular the inherited finite proof search remains unsuccessful even
after \emph{all} convergent Cayley relations and arbitrary shuffle products
of them have been included.
\end{theorem}

The distinction between the full ideal and its bounded intersection is
essential. We verify a finite basis in $F_4\mathcal A_7^-$, but the theorem
allows $\mathcal J$ in the entire admissible word algebra. The next two
subsections justify why those finite checks imply that stronger assertion.

\subsection{Why the depth-preserving projector controls the entire ideal}

We state and justify the precise prerequisite from \cite{gf:Depth}, so the
logic of the extension does not depend on a finite-rank extrapolation.
Introduce
\[
 x=c-b,\qquad y=a-\bar a,\qquad h=a+\bar a-b,
 \qquad \mathcal E=\{z,y,h,b,x\}.
\]
Then
\[
 c=x+b,\qquad a=(y+h+b)/2,\qquad
 \bar a=(-y+h+b)/2.
\]
In this basis admissibility means first letter different from $x$ and last
letter different from $z$. Depth is $|w|-n_z(w)$, conjugation acts by
$(-1)^{n_y(w)}$, and $C$ reverses a word, interchanges $z,x$, and
multiplies by $(-1)^{n_h(w)}$.

Let $\widetilde{\mathcal A}$ be the full shuffle algebra, including
nonadmissible formal words. It is used here only algebraically. With the
order $z<y<h<b<x$, every Lyndon word of length at least two is admissible.
The shuffle-polynomial theorem \cite{gf:Radford} consequently gives
\[
 \widetilde{\mathcal A}=\mathcal A[z,x].
\]
Let $R_0$ be the algebra retraction setting the two \emph{shuffle polynomial}
variables $z,x$ to zero. It preserves both letter counts, commutes with
$C,K$, and is explicitly
\begin{equation}
 R_0(w)=\sum_{u=0}^{l}\sum_{v=0}^{r}(-1)^{u+v}
          x^u\shuffle w[u:|w|-v]\shuffle z^v,
 \label{gf:eq:R0}
\end{equation}
where $l,r$ are the initial $x$ run and terminal $z$ run. Deleting an
initial $x$ and deleting a terminal $z$ are commuting shuffle derivations;
their finite Taylor substitutions prove \eqref{gf:eq:R0}. The factorials
cancel because shuffle powers of a single letter are factorial multiples
of its concatenation powers. In particular this is not literal erasure
of letters inside a word, and assigns no analytic value to divergent
integrals.

For nonempty $w$ put
\begin{equation}
 e(w)=\sum_{k=1}^{|w|}\frac{(-1)^{k-1}}k
       \sum_{w=w_1\cdots w_k}w_1\shuffle\cdots\shuffle w_k,
 \qquad E=R_0e,
 \label{gf:eq:E}
\end{equation}
where the blocks are consecutive and nonempty. The map $e$ is the first
Eulerian logarithm. It annihilates products of positive-degree factors and
equals the identity modulo such products \cite{gf:Patras}. These facts
also follow from the convolution logarithm of the identity character
into the commutative shuffle algebra: its logarithm is an infinitesimal
character, and the term $k=1$ gives the congruence. Thus $e^2=e$.
Reversal permutes consecutive cuts, so $e$ commutes with $C$; it also
commutes with $K$ and preserves both counts.

The subspace $U=E(\mathcal A)$ maps isomorphically onto the
indecomposables of $\mathcal A$, and multiplication identifies
$\mathcal A$ with $\operatorname{Sym}(U)$. To see this directly, $E$ kills
positive-degree products and is the identity modulo them. Any homogeneous
lift of an indecomposable basis therefore gives polynomial generators,
by a triangular change from the Lyndon generators. The lifts are bigraded
by $(n_z,n_x)$, and $C$ exchanges these two degrees.

For a homogeneous $u\in U$ of degree $(p,q)$ choose the orientation
\[
 Q_{\uparrow}u=
 \begin{cases}
 u,&p>q,\\
 Cu,&p<q,\\
 (u+Cu)/2,&p=q.
 \end{cases}
\]
On each pair of different bidegrees this identifies the two transported
generators. On the diagonal it retains the positive eigenspace of $C$ and
kills its negative eigenspace. Extending multiplicatively gives an
algebra retraction whose kernel is exactly $\mathcal J$: on the quotient
all $u-Cu$ vanish, hence $C$ acts identically on every polynomial, and
conversely $u-Cu$ are among the ideal generators. With
$E_{\uparrow}=Q_{\uparrow}E$, the retraction is the finite formula
\begin{equation}
 \Pi_{\uparrow}(w)=\sum_{k=1}^{|w|}\frac1{k!}
       \sum_{w=w_1\cdots w_k}
       E_{\uparrow}(w_1)\shuffle\cdots\shuffle E_{\uparrow}(w_k),
 \qquad \Pi_{\uparrow}(1)=1.
 \label{gf:eq:Pi}
\end{equation}
Convolution exponentiation of \eqref{gf:eq:E}, followed by $R_0$ and
the multiplicative orientation, proves this expression. For completeness,
$E(w)\in U$ even if a block $w$ is not admissible: in degree at least two,
the polynomial decomposition $\mathcal A[z,x]$ makes $w-R_0w$ a sum of
positive-degree products, killed by $e$; in degree one, $R_0e$ kills the
two extra generators directly.

The orientation never decreases the number of $z$'s in any nonzero
output. Shuffle multiplication adds their counts. Consequently, on
$\mathcal A$,
\begin{equation}
 \Pi_{\uparrow}^2=\Pi_{\uparrow},\quad
 \ker\Pi_{\uparrow}=\mathcal J,\quad
 \Pi_{\uparrow}K=K\Pi_{\uparrow},\quad
 \Pi_{\uparrow}(F_D\mathcal A)\subseteq F_D\mathcal A.
 \label{gf:eq:projector-properties}
\end{equation}
These are all-weight algebraic statements. Notice that
$C\Pi_{\uparrow}=\Pi_{\uparrow}$ is neither required nor generally true.

The verifier evaluates \eqref{gf:eq:E} using a permutation formula; its
equivalence with consecutive cuts is also elementary in every weight.
For a permutation of the $n$ positions let $d$ be the number of descents
of its inverse. A shuffle of consecutive blocks produces that permutation
exactly when cuts are placed at those $d$ required positions, with any
additional cuts chosen freely. Its total coefficient in \eqref{gf:eq:E}
is therefore
\[
 \sum_{k=d+1}^n\frac{(-1)^{k-1}}k
      \binom{n-1-d}{k-1-d}
   =(-1)^d\int_0^1 t^d(1-t)^{n-1-d}\,dt
   =\frac{(-1)^d}{n\binom{n-1}{d}}.
\]
This is precisely the implemented coefficient. Repeated letters simply
combine the integer multiplicities of equal output words. Thus the
finite audits of the two formulas are convention checks; their equality
at weight seven does not depend on extrapolating those audits.

\begin{lemma}[The full ideal intersection has a finite spanning certificate]
\label{gf:lem:intersection}
If $\mathcal B_4$ is any rational basis of $F_4\mathcal A_7^-$, then
\[
 \mathcal J_7^-\cap F_4\mathcal A_7
   =\operatorname{span}_{\mathbb Q}
       \{w-\Pi_{\uparrow}(w):w\in\mathcal B_4\}.
\]
\end{lemma}
\begin{proof}
Every displayed difference belongs to the full ideal by the kernel
property and belongs to the bounded space by depth preservation.
Conversely, for $f$ in the intersection, $\Pi_{\uparrow}f=0$, so
$f=(1-\Pi_{\uparrow})f$. Expanding $f$ in the finite basis proves the
opposite inclusion. No restriction on the construction of $f$ inside
$\mathcal J$ has been imposed.
\end{proof}

\subsection{The integer certificate and the proof}

Use the adapted-letter encoding $Z=0,Y=1,H=2,B=3,X=4$. Enumerate all
length-seven words lexicographically, retaining exactly those with first
letter different from $X$, last letter different from $Z$, at least three
$Z$'s, and an odd number of $Y$'s. This gives a basis $\mathcal B_4$ of
$F_4\mathcal A_7^-$ with $2546$ elements. The exact-depth counts are
$1,35,390,2120$. They can be checked directly from the enumeration, or
by choosing the nonzero positions and projecting to odd $Y$ parity.

The file \texttt{data/separator.json} supplies an integer functional
$\lambda$ on this basis. It is obtained by expressing the inherited
functional of \cite{gf:Bounded} in adapted coordinates; the functional
itself is not being claimed as a newly discovered vector. The new fact
is the second row of the following exact certificate:
\begin{align}
 \lambda(R)&=0 &&(R\in\mathscr R_4),\label{gf:eq:rows-zero}\\
 \lambda(w-\Pi_{\uparrow}w)&=0 &&(w\in\mathcal B_4),
       \label{gf:eq:full-ideal-zero}\\
 \lambda(T)&=
 -186660627289236812620583691130496682078843750753489297279472069700354208
 \ne0.\label{gf:eq:pairing}
\end{align}
The large integer is an exact value, not a numerical residual of the
polylogarithm conjecture.

\begin{proof}[Proof of Theorem~\ref{gf:thm:main}]
The standalone verifier regenerates the complete inherited row family,
checks \eqref{gf:eq:rows-zero}, and independently reconstructs each
$\Pi_{\uparrow}(w)$ by \eqref{gf:eq:R0}--\eqref{gf:eq:Pi}. It verifies
\eqref{gf:eq:full-ideal-zero} for all $2546$ basis words, and reconstructs
\eqref{gf:eq:target} before checking \eqref{gf:eq:pairing}. All arithmetic
is integral or rational. Thus Lemma~\ref{gf:lem:intersection} shows that
$\lambda$ annihilates $\mathcal J_7^-\cap F_4\mathcal A_7$.

If, contrary to the theorem, $T=j+r$ with $j\in\mathcal J_7^-$ and
$r\in\operatorname{span}\mathscr R_4$, then $T,r\in F_4\mathcal A_7^-$
implies $j=T-r\in F_4\mathcal A_7^-$. Consequently $\lambda(j)=0$;
also $\lambda(r)=0$ by \eqref{gf:eq:rows-zero}. This contradicts
$\lambda(T)\ne0$ and proves \eqref{gf:eq:nonmembership}.
\end{proof}

For reproducibility, execute
\begin{quote}\ttfamily
python code/verify\_complete\_cayley.py
\end{quote}
from the supplied contribution directory. It requires only Python's
standard library. The verifier checks SHA-256 digests of the two retained
source modules and the functional data before use. It regenerates $5131$
rows, verifies $2546$ complete-ideal basis differences, and checks the
nonzero target pairing. There are $2300$ nonzero projector differences;
no claim of their linear independence is made. Additional finite audits
compare the Eulerian permutation formula against consecutive cuts on all
$780$ words through weight four and check $2150$ shuffle products, as
well as retraction, conjugation, and depth conventions. These audits
support the implementation; the preceding algebraic proof supplies the
all-weight projector properties and the passage to the full ideal.

\subsection{Consequences for further research}

The immediate computational recommendation is sharper than in the
inspected sources. Completing Cayley symmetry while retaining precisely
$\mathscr R_4$ does not settle $S_6$. A successful continuation must change
the additional relation input: for example by admitting suitable
higher-depth double-shuffle or distribution rows, by adding another
functional equation, or by deriving a uniform generating identity that
implies the target. A larger numerical-precision calculation does not
address this missing algebraic information.

The theorem does not prove that every conceivable derivation must pass
through depth five. It excludes exactly the full ideal plus the specified
finite additional span; another depth-four family could carry information
outside that span. Likewise, no conclusion for the revised $S_8$ candidate
follows merely by analogy. Useful next questions are to identify the
lowest-weight additional row whose projected class pairs nontrivially
with $\lambda$, to test whether selected depth-five double-shuffle rows
suffice, and to find a conceptual functional equation explaining the
whole even-index family. These are exact proof-search questions, while
$R_6=0$ itself remains the separate numerical-period conjecture.
